% !TeX root = ../../../main.tex
\subsection{Background}
\label{ch02:sec:background}

We state the time series prediction setup we study in Section~\ref{ch02:sec:problem-setup}.
Our proposed method is based on supervised contrastive learning \citep{khosla2020supervised}, which we review in Section~\ref{ch02:sec:supcon}.
For any positive integer $k$, we regularly use the notation $[k]\triangleq \{1,2,\dots,k\}$.


\subsubsection{Problem Setup}
\label{ch02:sec:problem-setup}


% \vspace{.2em}
\noindent
\textbf{Training data.}
We assume that the training data consist of $N$ patients with different time series.
For the \mbox{$i$-th} patient (with $i\in[N]$), we observe $L_i$ time steps, where at each time step, we keep track of $D$ features (e.g., clinical measurements).
Specifically, we denote the \mbox{$i$-th} patient's feature vector at time step $\ell\in[L_i]$ (sorted chronologically) as $\mathbf{x}_i^{(\ell)}\in\mathbb{R}^{D}$. 
Moreover, we know the time-step times, where the \mbox{$i$-th} patient's time at the $\ell$-th time step is $t_i^{(\ell)}\in\mathbb{R}$.
This notation allows for features to be static, i.e., a feature could stay constant across time.

We further assume that every time step of a time series belongs to one of $C$ different classes, i.e., the set of classes is $[C]$.
In the dynamic outcome case, where the classification label varies over time, we assume that we know the classification label $y_i^{(\ell)}\in[C]$ for every patient
%data point
$i\in[N]$ for every time step $\ell\in[L_i]$.
For the static outcome case, we only know the classification label at the final time step per time series: we know $y_i^{(L_i)}\in[C]$ for all $i\in[N]$, and at earlier time steps $\ell < L_i$, we set $y_i^{(\ell)}=``?"$, which could be thought of as an additional ``unknown'' class.


\noindent
\textbf{Prediction task.} Given a test time series, suppose that we observe the feature vector $\mathbf{x}_*\in\mathbb{R}^{D}$ at a \emph{single time step}. We aim to predict the target label $y_*\in[C]$ corresponding to $\mathbf{x}_*$. In the dynamic outcome setting, this target label is for the same time step as $\mathbf{x}_*$. In the static outcome setting, this target label is for the label of final time step in the test time series.

We focus on predicting the classification outcome for each time step.
This captures the extreme case of a test time series consisting of a single time step. 
For example, this could happen in a clinical setup if a patient enters a new hospital for which we see measurements of this patient for the first time.
Our real-world clinical dataset experiments later reveal that our model outperforms baselines that leverage complete historical observations.

Importantly, although we focus on prediction given a single time step's feature vector $\mathbf{x}_*$, our problem setup allows for a feature in $\mathbf{x}_*$ to capture historical information (e.g., a feature in $\mathbf{x}_*$ could be the maximum observed white blood count over the past 4 hours).


\subsubsection{Supervised Contrastive Learning (SCL)}
\label{ch02:sec:supcon}


SCL learns an embedding representation of the data so that ``similar'' data points have embedding vectors that have higher cosine similarity compared to those of ``dissimilar'' data points.
Data points are ``similar'' if they have the same classification label.
As this framework was originally developed without temporal structure, we drop superscripts previously used to indicate dependence on time.


\noindent
\textbf{Notation.}
We use $f$ to denote the so-called \emph{encoder network}, where given any point~$\mathbf{x}$ in the raw input feature space, its embedding representation is $f(\mathbf{x})$.
This embedding representation is constrained to be a $d$-dimensional Euclidean vector with norm~1, also referred to as a \emph{hyperspherical embedding}.
Thus, $\|f(\mathbf{x})\|=1$.
We denote this hyperspherical output space as $\mathcal{S}^{d-1}\triangleq \{\mathbf{z}\in\mathbb{R}^d\text{ s.t.~}\|\mathbf{z}\|=1\}$.
For $\mathbf{u},\mathbf{v}\in\mathcal{S}^{d-1}$, the cosine similarity between $\mathbf{u}$ and $\mathbf{v}$ is $\frac{\langle \mathbf{u}, \mathbf{v} \rangle}{\|\mathbf{u}\|\|\mathbf{v}\|} = \langle \mathbf{u}, \mathbf{v} \rangle$, where $\langle\cdot,\cdot\rangle$ denotes the dot product.

We also denote $g$ as the \emph{predictor network} that maps the embedding output of the encoder in $\mathcal{S}^{d-1}$ to a probability distribution space over $C$ classes.


\noindent
\textbf{\bfseries\scshape Simple-SCL.}
For ease of exposition, we present a simplified version of SCL that we call \ChTwoSmallCaps{Simple-SCL}, \emph{which does not use data augmentation} (since our experiments later will be on tabular data without data augmentation as generating realistic fake patient data is challenging in the clinical setting).
We learn the encoder $f$ using minibatch gradient descent.
For a minibatch of $B$ training points $\mathbf{x}_1,\dots,\mathbf{x}_B$ with corresponding classification labels $y_1,\dots,y_B$, we denote the embedding vectors of these points as $\mathbf{z}_1=f(\mathbf{x}_1),\dots,\mathbf{z}_B=f(\mathbf{x}_B)$.
Since we want points with the same label to have high cosine similarity, we keep track of which points have the same label.
To do this, we let $\mathcal{P}(i)$ denote the set of points with the same label as the $i$-th point, excluding the $i$-th point:
\begin{equation}
\mathcal{P}(i) \triangleq  \{ j\in[B]\text{ s.t.~}y_j=y_i\text{ and }j\ne i \}.
\label{ch02:eq:neighbors-of-i}
\end{equation}
Next, we define the following ratio:
\begin{equation}
\Psi(i,j;\tau)
\triangleq  \frac{\exp(\langle \mathbf{z}_i, \mathbf{z}_j \rangle / \tau)}{\sum_{k\in[B]\text{ s.t.~}k\ne i}\exp(\langle \mathbf{z}_i, \mathbf{z}_k \rangle / \tau)},
\label{ch02:eq:Psi}
\end{equation}
where the constant $\tau>0$ is a user-specified hyperparameter. 
The key idea is that if the $i$-th and $j$-th points have the same label (so $j\in\mathcal{P}(i)$), then we want $\Psi(i,j;\tau)$ to be large: the numerator being large means that cosine similarity $\langle\mathbf{z}_i, \mathbf{z}_j\rangle$ is large while the denominator provides a normalization to ensure that $\Psi(i,j;\tau)\in[0,1]$.
Then to encourage $\Psi(i,j;\tau)$ to be large for all $i\in[B]$ and $j\in\mathcal{P}(i)$, we minimize the~loss
\begin{equation}
L_{\text{Simple-SCL}}
\triangleq 
-\!\!\!\!\!\!\!\!\!\!\!\!\!\!\!\sum_{i\in[B]\text{ s.t.~}|\mathcal{P}(i)|\ge1}
   \frac{1}{|\mathcal{P}(i)|}
     \sum_{j\in\mathcal{P}(i)} \log\Psi(i,j;\tau).
\label{ch02:eq:sup-con-loss}
\end{equation}
The original SCL uses a loss of the same form but with data augmentation (see Appendix~\ref{ch02:sec:relating-to-original-SCL}).

Lastly, in line with \cite{khosla2020supervised}, after the encoder is learned, we train a predictor network $g$ to map the resulting embeddings to a probability distribution over all classes.


